\section{Verwandte Arbeiten}

Die Einordnung in den wissenschaftlichen Kontext erfolgt entlang zweier Richtungen. Dies umfasst einerseits Arbeiten mit vergleichbarem energiewirtschaftlichen Anwendungsfokus sowie andererseits methodische Ansätze zur visuellen Verknüpfung mehrdimensionaler und dichter Zeitreihendarstellungen.

\subsection*{Visuelle Analyse von Energie-Zeitreihen}

Eine direkte Parallele zur vorliegenden Arbeit bildet die Untersuchung von van Wijk und van Selow \cite{vanWijk1999}. Die Autoren aggregieren univariate Zeitreihen des Energieverbrauchs mittels Clusteranalyse zu mittleren Tagesprofilen und visualisieren die Zuordnung der Einzeltage in einem kalendarischen Raster. Beiden Ansätzen liegt die Prämisse zugrunde, dass der Tag die zentrale Analyseeinheit für periodische Energiesysteme darstellt und tages- sowie jahreszeitliche Dynamiken simultan sichtbar sein müssen. Dieser Gedanke motiviert im vorliegenden Projekt das Tag-Stunden-Raster der Heatmap sowie die tagesbasierten Polygonzüge im Koordinatensystem. Wesentliche Unterschiede bestehen im Abstraktionsgrad und in der Dimensionalität. Während van Wijk und van Selow Cluster-Repräsentanten univariater Daten automatisiert berechnen, zeigt die hiesige Anwendung alle 364 Tage im Detail und überlässt die Gruppenbildung der interaktiven Filterung. Zudem verarbeitet die vorliegende Arbeit sieben gekoppelte Dimensionen je Tag, was den Einsatz paralleler Koordinaten begründet.

Janetzko et al. \cite{Janetzko2014} analysieren Stromverbrauchsdaten von Gebäuden mithilfe unüberwachter Anomalieerkennung, wobei errechnete Unregelmäßigkeiten über der Zeit visualisiert werden. Ähnlich wie im vorliegenden Bericht sollen seltene Extremzustände in hochaufgelösten Zeitreihen unter Berücksichtigung des jeweiligen Tagestyps hervorgehoben werden. Die Ansätze unterscheiden sich jedoch in der Aufgabenverteilung. Während bei Janetzko et al. ein statistischer Algorithmus Auffälligkeiten markiert, definieren Anwender im hiesigen System über Brushing flexibel eigene Kriterien für relevante Zustände. Dies gewährleistet eine transparente Nachvollziehbarkeit für energiewirtschaftliche Bewertungen. Inhaltlich verknüpft das vorliegende Projekt zudem die physikalische Erzeugung direkt mit dem Marktpreis, während sich Janetzko et al. auf reine Verbrauchsdaten beschränken.

\subsection*{Dichte Darstellungen, mehrdimensionale Profile und Kopplung}

Die theoretische Basis der pixelbasierten Heatmap geht auf Keim \cite{Keim2000} zurück, der den Entwurf dichter Darstellungen als Optimierungsaufgabe zur Erhaltung semantischer Distanzen beschreibt. Das Prinzip, jedem Messwert ein visuelles Element ohne vorherige Glättung zuzuordnen, wird in der hiesigen Anwendung aufgegriffen. Während Keim primär Daten ohne inhärente Ordnungsstruktur betrachtet und raumfüllende Spiral- oder Peano-Layouts entwirft, gibt die Zeitstruktur des Stromsystems eine natürliche zweidimensionale Matrix aus Tag und Stunde vor. Die methodische Grenze dichter Farbraster, wonach quantitative Werte nur näherungsweise aus Farbstufen abgelesen werden können, wird durch bedarfsgerechte Tooltips und Kennzahlenblöcke kompensiert.

Heinrich und Weiskopf \cite{Heinrich2013} bieten einen umfassenden Überblick über Verfahren für parallele Koordinaten. In Übereinstimmung mit ihrer Taxonomie dient die Technik hier vor allem dem interaktiven Filtern und dem Vergleich individueller Profile gegenüber der Hintergrundverteilung. Auf weitergehende Visualisierungsformen wie Kantenbündelung oder Dichteschätzungen wurde zugunsten einer klaren Lesbarkeit verzichtet. Die Achsenreihenfolge ist fachlich begründet fixiert, um die unmittelbare Interpretation ohne Neukonfiguration zu unterstützen.

Das Kopplungskonzept stützt sich auf die Systematisierung koordinierter Mehrfachansichten nach Roberts \cite{Roberts2007}. Zentrale Mechanismen wie verknüpfte Selektion, Hervorhebung und Filterung basieren hierbei auf einem gemeinsamen Anwendungszustand, der durch die funktionale Elm-Architektur deklarativ verwaltet wird. Während Roberts allgemeine Frameworks für frei verschaltbare Ansichten beschreibt, setzt diese Arbeit auf drei fest aufeinander abgestimmte Sichten, die dediziert aus den fachlichen Fragestellungen abgeleitet sind.

Aus beiden Forschungslinien ergibt sich die Ausrichtung dieser Arbeit. Sie übernimmt den Tag als fundamentale Analyseeinheit für Energiedaten, verzichtet jedoch zugunsten der Nachvollziehbarkeit auf vorgeschaltete automatisierte Clusteringverfahren. Die Kombination aus zweidimensionaler dichter Zeitdarstellung und mehrdimensionalen Koordinaten ermöglicht eine hypothesenfreie Exploration, deren Mehrwert in der unmittelbaren empirischen Überprüfbarkeit der Systemzusammenhänge liegt.
